\documentclass[12pt, a4paper]{article}

% ─── Packages (pdfLaTeX compatible) ──────────────────────────────────────────
\usepackage[utf8]{inputenc}
\usepackage[T5]{fontenc}
\usepackage[english]{babel}
\usepackage{geometry}
\usepackage{graphicx}
\usepackage{booktabs}
\usepackage{array}
\usepackage{tabularx}
\usepackage{amsmath, amssymb}
\usepackage{hyperref}
\usepackage{xcolor}
\usepackage{setspace}
\usepackage{titlesec}
\usepackage{fancyhdr}
\usepackage{caption}
\usepackage{subcaption}
\usepackage{float}
\usepackage{enumitem}
\usepackage{parskip}
\usepackage{tocloft}

% ─── Page Geometry ───────────────────────────────────────────────────────────
\geometry{
  a4paper,
  top    = 2.5cm,
  bottom = 2.5cm,
  left   = 3cm,
  right  = 2.5cm
}

% ─── Colours ─────────────────────────────────────────────────────────────────
\definecolor{usthblue}{RGB}{0, 70, 127}
\definecolor{usthred} {RGB}{196, 30, 58}

% ─── Hyperref Setup ──────────────────────────────────────────────────────────
\hypersetup{
  colorlinks = true,
  linkcolor  = black,
  urlcolor   = black,
  citecolor  = black,
  pdfauthor  = {Group Project},
  pdftitle   = {Assignment Report - Machine Learning and Data Mining II},
}

% ─── Header / Footer ─────────────────────────────────────────────────────────
\pagestyle{fancy}
\fancyhf{}
\fancyhead[L]{\small\textit{Machine Learning and Data Mining II}}
\fancyhead[R]{\small\textit{Book Recommendation System}}
\fancyfoot[C]{\thepage}
\renewcommand{\headrulewidth}{0.4pt}

% ─── Section Formatting ──────────────────────────────────────────────────────
\titleformat{\section}
  {\normalfont\large\bfseries}
  {\thesection.}{0.5em}{}
\titleformat{\subsection}
  {\normalfont\normalsize\bfseries}
  {\thesubsection}{0.5em}{}
\titleformat{\subsubsection}
  {\normalfont\normalsize\itshape}
  {\thesubsubsection}{0.5em}{}

% ─── ToC Formatting ──────────────────────────────────────────────────────────
\renewcommand{\cftsecleader}{\cftdotfill{\cftdotsep}}
\setlength{\cftbeforesecskip}{4pt}
\setlength{\cftbeforesubsecskip}{2pt}

% ─── Line Spacing ────────────────────────────────────────────────────────────
\setstretch{1.3}

% ═════════════════════════════════════════════════════════════════════════════
\begin{document}

% ─────────────────────────────────────────────────────────────────────────────
%  PAGE 1 – COVER PAGE
% ─────────────────────────────────────────────────────────────────────────────
\begin{titlepage}
  \thispagestyle{empty}
  \centering
  \setlength{\fboxrule}{1.5pt}
  \setlength{\fboxsep}{20pt}
  \fbox{%
    \begin{minipage}[c][0.90\textheight][c]{0.88\textwidth}
      \centering
      \vspace*{0.5cm}
      {\fontsize{13}{16}\selectfont\bfseries
        UNIVERSITY OF SCIENCE AND TECHNOLOGY OF HANOI\\[4pt]}
      {\fontsize{11}{14}\selectfont\itshape
        Faculty of Information and Communication Technology}
      \vspace{0.8cm}
      \IfFileExists{download.png}{%
        \includegraphics[width=5.5cm]{download.png}
      }{%
        \begingroup
          \setlength{\fboxrule}{0.6pt}
          \fbox{\parbox[c][3cm][c]{5.5cm}{\centering
            {\color{usthblue}\fontsize{28}{28}\selectfont\bfseries USTH}\\[4pt]
            {\color{usthred}\fontsize{8}{10}\selectfont VIETNAM FRANCE UNIVERSITY}
          }}
        \endgroup
      }
      \vspace{0.9cm}
      {\fontsize{12}{16}\selectfont\bfseries ASSIGNMENT REPORT\\}
      {\fontsize{12}{16}\selectfont for\\}
      {\fontsize{12}{16}\selectfont\bfseries\underline{MACHINE LEARNING AND DATA MINING II}\\}
      \vspace{0.7cm}
      {\fontsize{14}{18}\selectfont\bfseries Book Recommendation System\\[6pt]}
      {\fontsize{11}{14}\selectfont\bfseries Dr. Doan Nhat Quang\\}
      \vspace{0.9cm}
      \begin{tabular}{|>{\centering\arraybackslash}m{1.2cm}
                       |>{\centering\arraybackslash}m{5.5cm}
                       |>{\centering\arraybackslash}m{3cm}|}
        \hline
        \textbf{No.} & \textbf{Name} & \textbf{Student ID} \\
        \hline
        1 & Nguyễn Duy Hưng & 22BA13154 \\
        \hline
        2 & Nguyễn Tiến Dũng & 22BA13087 \\
        \hline
        3 & Đàm Tuấn Hưng & 22BA13150 \\
        \hline
        4 & Lại Quang Huy & 22BA13161 \\
        \hline
        5 & Hoàng Nhật Anh & 22BA13011 \\
        \hline
      \end{tabular}
      \vspace{0.5cm}
    \end{minipage}%
  }
\end{titlepage}

% ─────────────────────────────────────────────────────────────────────────────
%  PAGE 2 – LIST OF ABBREVIATIONS
% ─────────────────────────────────────────────────────────────────────────────
\newpage
\thispagestyle{fancy}
\section*{List of Abbreviations}
\addcontentsline{toc}{section}{List of Abbreviations}

\begin{tabular}{@{}ll@{}}
  \textbf{CB}   & Content-Based (Filtering) \\
  \textbf{CF}   & Collaborative Filtering \\
  \textbf{EDA}  & Exploratory Data Analysis \\
  \textbf{KNN}  & K-Nearest Neighbors \\
  \textbf{LOO}  & Leave-One-Out \\
  \textbf{MAE}  & Mean Absolute Error \\
  \textbf{ML}   & Machine Learning \\
  \textbf{RMSE} & Root Mean Square Error \\
  \textbf{RS}   & Recommender System \\
  \textbf{SVD}  & Singular Value Decomposition \\
  \textbf{TF-IDF} & Term Frequency -- Inverse Document Frequency \\
\end{tabular}

% ─────────────────────────────────────────────────────────────────────────────
%  PAGE 3 – ABSTRACT
% ─────────────────────────────────────────────────────────────────────────────
\newpage
\thispagestyle{fancy}
\section*{Abstract}
\addcontentsline{toc}{section}{Abstract}

The rapid growth of online book platforms has created an information overload problem,
making it difficult for readers to discover books that match their preferences.
This report presents a Book Recommendation System developed in two complementary parts,
both evaluated on the GoodBooks-10K dataset.

\textbf{Part~1 (Collaborative Filtering Baseline)} implements and compares two pure CF approaches:
\textbf{KNN-based CF} (memory-based, User-User and Item-Item variants) and
\textbf{SVD-based CF} (model-based matrix factorization trained with PyTorch).
Both methods rely solely on rating history, which limits their ability to serve
new users with no prior interactions --- the well-known \textit{Cold Start problem}.

\textbf{Part~2 (Hybrid Extension)} addresses this limitation by combining
Content-Based KNN with Truncated SVD collaborative signals.
The Content-Based component represents books as TF-IDF feature vectors and
uses cosine-distance KNN to recommend books from item metadata alone,
requiring zero user history.
For returning users, the system blends content and collaborative signals in a
weighted combination to produce personalised recommendations.
An additional smart \textit{intent detection} mechanism supports both
exact title lookup and free-text semantic discovery within a single interface.

Experimental results show that the Hybrid Extension significantly outperforms
the CF baseline in terms of Cold Start handling and recommendation relevance,
while the CF baseline provides stronger rating prediction accuracy.
Both approaches are compared and analysed to highlight their respective trade-offs.

% ─────────────────────────────────────────────────────────────────────────────
%  TABLE OF CONTENTS
% ─────────────────────────────────────────────────────────────────────────────
\newpage
\thispagestyle{fancy}
\tableofcontents

% ─────────────────────────────────────────────────────────────────────────────
%  LIST OF FIGURES AND TABLES
% ─────────────────────────────────────────────────────────────────────────────
\newpage
\thispagestyle{fancy}
\listoffigures
\addcontentsline{toc}{section}{List of Figures}
\listoftables
\addcontentsline{toc}{section}{List of Tables}

% ─────────────────────────────────────────────────────────────────────────────
%  SECTION 1 – INTRODUCTION
% ─────────────────────────────────────────────────────────────────────────────
\newpage
\section{Introduction}

\subsection{Background and Motivation}

In recent years, the internet has changed the way people read books.
Online platforms such as Amazon, Goodreads, and Google Books now offer millions of book titles to users around the world.
Although this gives readers many choices, it also causes a problem: people often spend a lot of time searching for a book they like but cannot find the right one easily.
This situation is sometimes called \textit{information overload} \cite{resnick1994}.

To help solve this problem, recommender systems have been developed.
A recommender system is a tool that suggests items to users based on their past behaviour or preferences.
In the context of books, such a system can look at a user's reading history and ratings to suggest books that the user might enjoy \cite{ricci2011}.

In this project, we use \textit{Collaborative Filtering} (CF) as the main approach for Part~1.
CF is based on the idea that people who liked the same books in the past are likely to have similar tastes in the future \cite{koren2009}.
Within the CF framework, we explore two methods: \textbf{KNN-based CF} (memory-based) and \textbf{SVD-based CF} (model-based matrix factorization) \cite{ricci2011}.
Part~2 extends this with a Hybrid approach that incorporates Content-Based filtering
to address the Cold Start limitation identified in Part~1.

\subsection{Problem Statement}

Although recommender systems have been studied extensively, several challenges remain.
First, book rating datasets are usually very \textit{sparse}: most users only rate a small number of books,
making it harder to compute reliable similarity scores \cite{su2009}.
Second, the \textit{cold start problem} arises when a new user joins the platform,
as there is no rating history to base recommendations on \cite{schein2002}.
Third, within CF, key design choices such as the number of neighbours $k$ for KNN
and the number of latent factors for SVD must be carefully tuned.
This project addresses these issues by implementing, comparing, and extending both approaches.

\subsection{Objectives}

The main goals of this project are:

\begin{itemize}[leftmargin=2cm]
  \item To collect and preprocess the GoodBooks-10K dataset (10,000 books, 5,976,479 ratings).
  \item To implement and compare two pure CF methods: KNN-based CF and SVD-based CF.
  \item To identify the Cold Start weakness of pure CF and propose a Hybrid extension.
  \item To build a Hybrid Content-Based KNN + TruncatedSVD system that serves both new and returning users.
  \item To evaluate all approaches, quantify the improvement, and deploy a web application.
\end{itemize}

The rest of this report is structured as follows.
Section~\ref{sec:dataset} explains the dataset and preprocessing.
Section~\ref{sec:methodology} describes the CF methods used in Part~1.
Section~\ref{sec:results} presents the CF baseline results and evaluation metrics.
Section~\ref{sec:hybrid} introduces the Hybrid Extension (Part~2) and its evaluation.
Section~\ref{sec:discussion} discusses findings, and Section~\ref{sec:conclusion} concludes.

% ─────────────────────────────────────────────────────────────────────────────
%  SECTION 2 – DATASET AND PREPROCESSING
% ─────────────────────────────────────────────────────────────────────────────
\newpage
\section{Dataset and Preprocessing}
\label{sec:dataset}

\subsection{Dataset Description}

The dataset used in this project is the \textbf{GoodBooks-10K} dataset,
containing real book ratings collected from the Goodreads platform.
It consists of two main CSV files: \texttt{goodbooks\_ratings.csv} and \texttt{goodbooks\_books.csv}.

\textbf{File 1: \texttt{goodbooks\_ratings.csv}} stores rating interactions between users and books.
Table~\ref{tab:ratings_schema} shows its schema.

\begin{table}[H]
  \centering
  \caption{Schema of \texttt{goodbooks\_ratings.csv}}
  \label{tab:ratings_schema}
  \begin{tabular}{|l|l|l|}
    \hline
    \textbf{Field}   & \textbf{Type} & \textbf{Description} \\
    \hline
    \texttt{user\_id}  & int & User ID (1 to 53,424) \\
    \hline
    \texttt{book\_id}  & int & Book ID (1 to 10,000) \\
    \hline
    \texttt{rating}    & int & Rating score: 1 to 5 stars \\
    \hline
  \end{tabular}
\end{table}

Key statistics: 5,976,479 total ratings; 53,424 users; 10,000 books.
Ratings per user average $\approx$111.9 (median 70);
ratings per book average $\approx$597.6 (median 430).
The user-item matrix has an initial sparsity of $\approx$98.88\%.

\textbf{File 2: \texttt{goodbooks\_books.csv}} contains metadata for 10,000 books,
including \texttt{book\_id}, \texttt{title}, \texttt{authors}, \texttt{average\_rating},
\texttt{description}, \texttt{genres}, and \texttt{num\_pages}.
This file is used by the Content-Based component in Part~2.

\subsection{Data Exploration}

Figure~\ref{fig:eda} shows four EDA charts: the overall rating distribution,
ratings per user (capped at 300), ratings per book (capped at 2000),
and the Top-10 most-rated books.

\begin{figure}[H]
  \centering
  \includegraphics[width=\linewidth]{eda_overview.png}
  \caption{EDA overview --- rating distribution, user activity, book popularity, top-10 books.}
  \label{fig:eda}
\end{figure}

Two key observations emerge from the EDA.
First, ratings are positively skewed (4--5 stars dominate), reflecting the fact that
users tend to rate books they have already chosen to read and enjoyed.
Second, book popularity follows a long-tail distribution: a small number of books
attract the vast majority of ratings, while most books receive very few.
These two characteristics directly motivate the filtering thresholds applied in Section~\ref{sec:cleaning}.

Figure~\ref{fig:sparsity} visualises the user-item matrix for the first 150 users and 150 books,
confirming the 98.88\% sparsity level that motivates our modelling choices.

\begin{figure}[H]
  \centering
  \includegraphics[width=0.85\linewidth]{sparsity_heatmap.png}
  \caption{User-Item Matrix Sparsity Heatmap (first 150 users $\times$ 150 books).}
  \label{fig:sparsity}
\end{figure}

\subsection{Data Cleaning}
\label{sec:cleaning}

To ensure that similarity calculations are reliable, we remove users and books
that do not meet minimum rating thresholds.
This step reduces noise, slightly increases matrix density, and ensures that
every user retained in the dataset can always receive recommendations.

The filtering was done in two passes, as shown in Table~\ref{tab:filtering}:

\begin{table}[H]
  \centering
  \caption{Data filtering thresholds}
  \label{tab:filtering}
  \begin{tabular}{|l|c|p{7cm}|}
    \hline
    \textbf{Condition}       & \textbf{Threshold} & \textbf{Reason} \\
    \hline
    Ratings per user $\geq$  & 20 & Remove inactive users without enough history for CF \\
    \hline
    Ratings per book $\geq$  & 50 & Remove rarely-rated books with unreliable similarity \\
    \hline
  \end{tabular}
\end{table}

After filtering, 8,000 users were randomly sampled to keep the matrix size within memory limits.
The resulting matrix for the CF baseline has \textbf{8,000 users $\times$ 9,991 books}.
Note that the Hybrid Extension (Part~2) uses all 53,424 users for SVD training,
as TruncatedSVD operates on the full sparse matrix without memory issues.

\subsection{Train-Test Split}

\subsubsection*{Leave-One-Out (LOO) Method}

For the train-test split, we use the \textbf{Leave-One-Out (LOO)} strategy.
For each user, the book with the highest rating is held out as the test item;
all other ratings form the training set.
This approach simulates real-world next-item prediction, ensures every user appears
in both train and test sets, and enables evaluation via Top-$N$ Hit Rate metrics.

\subsubsection*{Mean-Centering (KNN only)}

Before computing similarity for the KNN model, ratings are \textit{mean-centered} per user:
\[
  r'(u, i) = r(u, i) - \bar{r}(u)
\]
This removes user rating bias (e.g., a strict user who always rates lower).
For SVD, user and item biases are handled automatically through dedicated bias terms
$b_u$ and $b_i$ learned during training.

% ─────────────────────────────────────────────────────────────────────────────
%  SECTION 3 – METHODOLOGY
% ─────────────────────────────────────────────────────────────────────────────
\newpage
\section{Methodology}
\label{sec:methodology}

\subsection{Problem Formulation}

Let $U$ be the set of users and $I$ be the set of books.
Each user $u \in U$ has rated a subset $I_u \subseteq I$ with ratings $r(u, i) \in \{1,2,3,4,5\}$.
The task is to predict $\hat{r}(u, i)$ for all unrated books $i \notin I_u$,
and return the Top-10 books with the highest predicted scores.

\subsection{Collaborative Filtering Framework}

Collaborative Filtering (CF) relies solely on user rating history to find patterns \cite{su2009}.
We implement CF using two different methods:
\textbf{KNN-based CF} (memory-based, using user-user or item-item similarity) and
\textbf{SVD-based CF} (model-based, using PyTorch neural matrix factorization).
Both are applied to the same preprocessed dataset under the same conditions.

\subsection{KNN-based Collaborative Filtering}

\subsubsection{User-Based Collaborative Filtering (User-User KNN)}

\textbf{Objective:} Recommend books based on what users similar to $u$ have rated highly.

\textbf{Step 1 --- Compute User-User Similarity.}
Cosine similarity is applied on the mean-centered rating matrix $R'$:
\[
  \text{sim}(u, v) = \frac{\mathbf{r'}_u \cdot \mathbf{r'}_v}{\|\mathbf{r'}_u\| \cdot \|\mathbf{r'}_v\|}
\]

\textbf{Step 2 --- Select $K$ Nearest Neighbors.}
For predicting user $u$'s rating for book $i$, we select $K = 20$ users
with the highest similarity to $u$ who have also rated book $i$.

\textbf{Step 3 --- Predict Rating.}
\[
  \hat{r}(u, i) = \bar{r}(u) + \frac{\sum_{v \in \mathcal{N}(u)} \text{sim}(u,v) \cdot r'(v,i)}
                                     {\sum_{v \in \mathcal{N}(u)} |\text{sim}(u,v)|}
\]

\textbf{Step 4 --- Generate Top-$N$ List.}
Predicted scores are accumulated across all $K$ neighbors for all unread books;
the Top-10 with the highest scores are returned.

\subsubsection{Item-Based Collaborative Filtering (Item-Item KNN)}

\textbf{Objective:} Recommend books similar to those the user has already rated highly.

\textbf{Step 1 --- Compute Item-Item Similarity.}
The rating matrix is transposed to books $\times$ users format, then cosine similarity is applied:
\[
  \text{sim}(i, j) = \frac{\mathbf{c}_i \cdot \mathbf{c}_j}{\|\mathbf{c}_i\| \cdot \|\mathbf{c}_j\|}
\]

\textbf{Step 2 --- Predict Rating.}
For user $u$ predicting book $i$, we find $K$ books rated by $u$ most similar to $i$:
\[
  \hat{r}(u, i) = \bar{r}(u) + \frac{\sum_{j \in \mathcal{N}(i)} \text{sim}(i,j) \cdot r'(u,j)}
                                     {\sum_{j \in \mathcal{N}(i)} |\text{sim}(i,j)|}
\]

\textbf{Step 3 --- Generate Top-$N$ List.}
For each book the user has read, similarity scores are accumulated for unread books;
the Top-10 with the highest total scores are returned.

\subsection{SVD-based Collaborative Filtering \cite{funk2006, koren2009}}
\label{sec:svd_pytorch}

\subsubsection{Model Architecture}

The predicted rating is computed as:
\[
  \hat{r}(u, i) = \mu + b_u + b_i + {q}_i^\top {p}_u
\]
where the components are defined in Table~\ref{tab:svd_components}.

\begin{table}[H]
  \centering
  \caption{Components of the SVD model}
  \label{tab:svd_components}
  \begin{tabular}{|l|c|p{7.5cm}|}
    \hline
    \textbf{Component}     & \textbf{Symbol}              & \textbf{Meaning} \\
    \hline
    Global Bias            & $\mu$                        & Overall platform rating tendency ($\approx$3.85 stars) \\
    \hline
    User Bias              & $b_u$                        & Personal tendency (e.g., a strict user always rates lower) \\
    \hline
    Item Bias              & $b_i$                        & Book-level tendency (popular books tend to be rated higher) \\
    \hline
    Latent Interaction     & ${q}_i^\top {p}_u$ & Dot product capturing alignment of user and book factors \\
    \hline
  \end{tabular}
\end{table}

This is implemented using \textbf{PyTorch} with user and book embeddings as \texttt{nn.Embedding} layers.
The prediction separates into a \textit{bias} part ($\mu + b_u + b_i$) capturing general tendencies,
and an \textit{interaction} part (${q}_i^\top {p}_u$) capturing the personal match between user and book.

\subsubsection{Training Process}

\begin{enumerate}
  \item \textbf{Continuous ID Mapping.} User and book IDs are remapped to 0-based indices.
  \item \textbf{Mini-Batch DataLoader.} 5,976,479 ratings are split into batches of 10,000.
  \item \textbf{Weight Initialisation.}
        Latent vectors ${p}_u, {q}_i \sim \mathcal{N}(0, 0.01)$;
        biases initialised to 0; $\mu$ set to the mean of all training ratings.
  \item \textbf{Training Loop (10 Epochs).}
        Forward pass $\rightarrow$ MSE Loss $\rightarrow$
        Backpropagation $\rightarrow$ Adam update (lr = 0.01).
  \item \textbf{Inference.}
        All unread books for user $u$ are scored, sorted, and the Top-10 returned.
\end{enumerate}

Table~\ref{tab:svd_convergence} confirms steady convergence over 10 epochs.

\begin{table}[H]
  \centering
  \caption{Training convergence of the SVD model}
  \label{tab:svd_convergence}
  \begin{tabular}{|c|c|}
    \hline
    \textbf{Epoch} & \textbf{Train RMSE} \\
    \hline
    1  & $\approx$ 0.88 \\
    \hline
    5  & $\approx$ 0.54 \\
    \hline
    10 & $\approx$ 0.45 \\
    \hline
  \end{tabular}
\end{table}

% ─────────────────────────────────────────────────────────────────────────────
%  SECTION 4 – EXPERIMENTAL RESULTS
% ─────────────────────────────────────────────────────────────────────────────
\newpage
\section{Experimental Results and Comparison}
\label{sec:results}

\subsection{Evaluation Metrics}

We use two groups of metrics to evaluate the CF models comprehensively.

\textbf{Group 1 --- Rating Prediction.}

\begin{table}[H]
  \centering
  \caption{Rating prediction metrics}
  \label{tab:metrics_rating}
  \begin{tabular}{|l|l|p{6.5cm}|}
    \hline
    \textbf{Metric} & \textbf{Formula} & \textbf{Meaning} \\
    \hline
    RMSE & $\sqrt{\frac{1}{n}\sum({\hat{r} - r})^2}$ & Sensitive to large errors; unit is stars \\
    \hline
    MAE  & $\frac{1}{n}\sum|\hat{r} - r|$             & Average absolute deviation; easier to interpret \\
    \hline
  \end{tabular}
\end{table}

\textbf{Group 2 --- Top-$N$ Recommendation Quality (Leave-One-Out).}
For each user, one test item is hidden; the model generates a Top-10 list.
A ``hit'' occurs when the hidden item appears in the Top-10 list.

\begin{table}[H]
  \centering
  \caption{Top-$N$ recommendation quality metrics (LOO)}
  \label{tab:metrics_topn}
  \begin{tabular}{|l|l|p{6cm}|}
    \hline
    \textbf{Metric}    & \textbf{Formula}           & \textbf{Meaning} \\
    \hline
    HitRate@10         & \# hits / total users       & Fraction of users with a relevant item in Top-10 \\
    \hline
    Precision@10       & \# hits / 10                & Fraction of relevant items in the Top-10 list \\
    \hline
    Recall@10          & \# hits / 1                 & Coverage (equals HitRate since each user has 1 test item) \\
    \hline
  \end{tabular}
\end{table}

\textit{Note:} Under LOO, Recall@10 $=$ HitRate@10 by definition.

\subsection{Experimental Setup}

All three CF models were trained and evaluated on the same preprocessed dataset
(8,000 users $\times$ 9,991 books, LOO split) under the conditions in Table~\ref{tab:config}.

\begin{table}[H]
  \centering
  \caption{Common experimental configuration (CF Baseline)}
  \label{tab:config}
  \begin{tabular}{|c|c|}
    \hline
    \textbf{Parameter}               & \textbf{Value} \\
    \hline
    Sampled users (after filtering)  & 8,000 \\
    \hline
    Books (after filtering)          & 9,991 \\
    \hline
    Split method                     & Leave-One-Out (LOO) \\
    \hline
    Users used for evaluation        & 500 \\
    \hline
    Top-$N$                          & 10 \\
    \hline
    $K$ Neighbors (KNN)              & 20 \\
    \hline
    Similarity metric (KNN)          & Cosine Similarity \\
    \hline
    Latent Factors (SVD)             & 100 \\
    \hline
    Epochs (SVD)                     & 10 \\
    \hline
    Optimizer (SVD)                  & Adam (lr = 0.01) \\
    \hline
    Batch Size (SVD)                 & 10,000 \\
    \hline
  \end{tabular}
\end{table}

\subsection{Results of KNN-based CF}

Figure~\ref{fig:rating_pred} shows RMSE and MAE for both KNN variants.

\begin{figure}[H]
  \centering
  \includegraphics[width=0.85\linewidth]{rating_prediction.png}
  \caption{Rating Prediction Error --- RMSE vs MAE (CF baseline)}
  \label{fig:rating_pred}
\end{figure}

Both variants achieve RMSE $<$ 1, indicating reasonable rating prediction accuracy.
Item-Based KNN performs more consistently (RMSE = 0.889, MAE = 0.669) compared to
User-Based KNN (RMSE = 0.908, MAE = 0.696).
This is likely because book profiles are more stable than user profiles:
a book's content and its rater group do not change over time,
making item-item similarity more reliable than user-user similarity.

\subsection{Results of SVD-based CF}

Figure~\ref{fig:svd_results} shows how Precision@K, Recall@K, and F1@K change
as the recommendation list size $K$ increases from 5 to 20.

\begin{figure}[H]
  \centering
  \includegraphics[width=0.85\linewidth]{svd_results.png}
  \caption{SVD --- Precision, Recall \& F1 @ K, threshold $\geq$ 4 (CF baseline)}
  \label{fig:svd_results}
\end{figure}

As expected, Precision@K decreases slightly as $K$ grows (longer lists include less-relevant items),
while Recall@K increases (more relevant items are covered).
F1@K remains stable and improves gradually, suggesting the model maintains
a good balance between accuracy and coverage.
A recommendation list size of $K \approx 15$--$20$ offers the best F1 trade-off for SVD.

\subsection{Model Comparison}

Table~\ref{tab:results_compare} summarises performance across all five metrics.
Item-Based KNN achieves the best scores on every metric.

\begin{table}[H]
  \centering
  \caption{Results comparison across all CF models}
  \label{tab:results_compare}
  \begin{tabular}{|l|c|c|c|c|c|}
    \hline
    \textbf{Model} & \textbf{Precision@10} & \textbf{Recall@10} & \textbf{HitRate@10} & \textbf{RMSE} & \textbf{MAE} \\
    \hline
    User-Based KNN  & 0.0128 & 0.1280 & 0.1280 & 1.0446 & 0.8959 \\
    \hline
    \textbf{Item-Based KNN} & \textbf{0.0200} & \textbf{0.2000} & \textbf{0.2000} & \textbf{0.9811} & \textbf{0.8292} \\
    \hline
    SVD             & 0.0014 & 0.0140 & 0.0140 & 1.1537 & 0.8799 \\
    \hline
  \end{tabular}
\end{table}

Figure~\ref{fig:model_comparison} visualises these results,
with the left panel comparing Top-10 ranking metrics and
the right panel comparing rating prediction error.

\begin{figure}[H]
  \centering
  \includegraphics[width=\linewidth]{model_comparison.png}
  \caption{Collaborative Filtering --- All Models Comparison (CF baseline)}
  \label{fig:model_comparison}
\end{figure}

Three conclusions can be drawn from these results.
First, Item-Based KNN outperforms all other models on every metric,
benefiting from stable book profiles and a more tractable item-item matrix
(9,991 books vs.\ 8,000 users).
Second, SVD achieves a low HitRate@10 (1.4\%) because it is optimised for
rating prediction (RMSE minimisation) rather than ranking.
With only 10 training epochs and no regularisation, the model has not fully converged;
in practice, SVD-based models perform strongly with more training and BPR loss \cite{koren2009}.
Third, HitRate@10 = 20\% represents a meaningful result: with 9,991 books,
a random model would achieve approximately 0.1\%, making Item-Based KNN
roughly 200$\times$ better than random \cite{su2009}.

% ─────────────────────────────────────────────────────────────────────────────
%  SECTION 5 – HYBRID EXTENSION
% ─────────────────────────────────────────────────────────────────────────────
\newpage
\section{Extension: Hybrid Content-Based + Collaborative Filtering System}
\label{sec:hybrid}

\subsection{Motivation: Cold Start in Pure Collaborative Filtering}

A critical finding from Section~\ref{sec:results} is that all three CF models share a
fundamental weakness: they require a user's rating history to generate recommendations.
This is the \textbf{Cold Start problem} \cite{schein2002}.
In practice, a significant proportion of platform users have rated fewer than 5 books,
making CF unreliable or impossible for them.
Table~\ref{tab:coldstart_gap} quantifies this gap.

\begin{table}[H]
  \centering
  \caption{Cold Start failure analysis on GoodBooks-10K}
  \label{tab:coldstart_gap}
  \begin{tabular}{|l|c|c|}
    \hline
    \textbf{Scenario} & \textbf{Pure CF (baseline)} & \textbf{Hybrid Extension} \\
    \hline
    New user, no rating history & \textbf{Fails} & \textbf{Works} (Content-Based KNN) \\
    \hline
    User with $< 20$ ratings & Falls back to global average & KNN (content) $+$ partial SVD \\
    \hline
    Active user ($\geq 20$ ratings) & Works normally & KNN $+$ SVD blended \\
    \hline
    Free-text query (no title) & Not supported & TF-IDF semantic discovery \\
    \hline
  \end{tabular}
\end{table}

\subsection{Hybrid System Architecture}

% [NOTE] This section uses TruncatedSVD from scikit-learn for collaborative
% signals --- a matrix decomposition method distinct from the PyTorch neural
% SVD trained in Part 1 (Section~\ref{sec:svd_pytorch}). While both perform
% Singular Value Decomposition, TruncatedSVD operates directly on the sparse
% user-item matrix without gradient descent, making it faster to apply to
% the full 53,424-user matrix but without learned bias terms.

The Hybrid system is composed of three cooperating components:

\begin{enumerate}
  \item \textbf{Content-Based KNN} (\texttt{src/knn\_model.py}).
        Books are represented as TF-IDF feature vectors.
        Given a seed book, cosine-distance KNN finds the 10 most similar books
        from metadata alone --- no user history required.

  \item \textbf{TruncatedSVD} (\texttt{src/svd\_model.py}).
        The user-item rating matrix $R \in \mathbb{R}^{53{,}424 \times 10{,}000}$
        is decomposed retaining $k=50$ latent factors:
        \[
          R \approx U \Sigma V^\top
        \]
        The reconstructed matrix $\hat{R} = U \Sigma V^\top$ gives a predicted score
        for any (user, book) pair, capturing collaborative taste signals.

  \item \textbf{Hybrid Router} (\texttt{src/hybrid.py}).
        Selects the appropriate strategy based on user context (see Section~\ref{sec:routing}).
\end{enumerate}

\subsection{Feature Engineering for Content-Based KNN}

A \textit{text soup} is constructed for each book by concatenating fields,
with genre and author repeated twice to increase their relative weight:
\[
  \text{soup}_i = \text{title}_i \;\|\;
                  2\times\text{authors}_i \;\|\;
                  2\times\text{genres}_i \;\|\;
                  \text{description}_i
\]

This soup is vectorised by \texttt{TfidfVectorizer(max\_features=5000, stop\_words=`english')}
into a sparse matrix $T \in \mathbb{R}^{10000 \times 5000}$.
Two normalised numerical features (average\_rating, page count) are appended:
\[
  \mathbf{M} = [\,T \;|\; N\,] \in \mathbb{R}^{10000 \times 5003}
\]

The KNN model is trained as
\texttt{NearestNeighbors(n\_neighbors=11, metric=`cosine', algorithm=`brute')} on $\mathbf{M}$.
We request 11 neighbors and skip index~0 (the query book itself), returning 10 recommendations.

\begin{table}[H]
  \centering
  \caption{Content-Based KNN configuration}
  \label{tab:knn_config}
  \begin{tabular}{|l|c|p{6cm}|}
    \hline
    \textbf{Parameter} & \textbf{Value} & \textbf{Rationale} \\
    \hline
    TF-IDF max\_features & 5,000 & Captures vocabulary without overfitting \\
    \hline
    Numerical features & 3 & avg\_rating, ratings\_count, pages \\
    \hline
    Feature dimensions & 5,003 & Combined TF-IDF + numerical vector \\
    \hline
    Similarity metric & Cosine & Scale-invariant; optimal for TF-IDF space \\
    \hline
    $n\_neighbors$ & 11 & Return top-10; skip self at index 0 \\
    \hline
  \end{tabular}
\end{table}

\subsection{Hybrid Routing Logic}
\label{sec:routing}

The \texttt{HybridRecommender} routes each query through one of three strategies:

\begin{enumerate}
  \item \textbf{Cold Start} --- \texttt{user\_id} not provided:
        Return pure KNN recommendations based on content similarity.
  \item \textbf{Light User} --- \texttt{user\_id} provided but fewer than 20 ratings:
        KNN fallback, since TruncatedSVD latent factors are unreliable for sparse users.
  \item \textbf{Active User} --- \texttt{user\_id} with $\geq 20$ ratings:
        \[
          \text{final\_score} = 0.4 \times \text{score}_{\text{KNN}} + 0.6 \times \text{score}_{\text{SVD}}
        \]
        SVD weight is higher (60\%) because personalised signals dominate content similarity
        once sufficient interaction data is available.
\end{enumerate}

\subsection{Smart Intent Detection}

The web interface adds a single-input smart search that automatically detects intent:

\begin{itemize}
  \item \textbf{Search Mode} --- the query exactly matches a title in the dataset.
        The system displays a ``Book Detail'' card with metadata, then recommends similar books.
  \item \textbf{Discovery Mode} --- no exact title match found.
        The query is embedded as a TF-IDF vector and matched against all 10,000 books
        by cosine similarity, returning the most semantically relevant titles.
\end{itemize}

\subsection{Dataset Overview --- Hybrid System}

Figure~\ref{fig:dataset_overview} shows the dataset characteristics used by the Hybrid system.
All 5,976,479 ratings, 10,000 books, and 53,424 users are used for TruncatedSVD;
content features are built from all 10,000 books.

\begin{figure}[H]
  \centering
  \includegraphics[width=\linewidth]{1_dataset_overview.png}
  \caption{Dataset Overview --- GoodBooks-10K: rating distribution, top genres,
           ratings-per-user distribution, and book average rating distribution.}
  \label{fig:dataset_overview}
\end{figure}

Key observations: ratings are positively skewed (4--5 stars dominate);
Fiction, Fantasy, and Romance are the three most common genres;
median ratings per user = 111; book average ratings cluster around 3.5--4.5 (mean: 4.00).

\subsection{Evaluation --- Content-Based KNN Precision@K}

Figure~\ref{fig:precision_k} shows the Precision@K curve evaluated over 60 active users
using a Leave-One-Out approach: one liked book (rating $\geq$ 4) serves as the seed;
all other liked books form the ground truth.

\[
  \text{Precision@K} = \frac{|\text{Top-}K_{\text{rec}} \cap \text{Ground Truth}|}{K}
\]

\begin{figure}[H]
  \centering
  \includegraphics[width=0.85\linewidth]{2_knn_precision_at_k.png}
  \caption{Content-Based KNN --- Precision@K curve (K $= 1, 3, 5, 10, 15, 20$)}
  \label{fig:precision_k}
\end{figure}

\begin{table}[H]
  \centering
  \caption{Content-Based KNN Precision@K results}
  \label{tab:cb_precision}
  \begin{tabular}{|c|c|c|}
    \hline
    $K$ & Precision@K & Improvement vs.~Random (1 in 10k) \\
    \hline
    1  & 15.0\% & $1500\times$ \\
    3  & 15.0\% & $1500\times$ \\
    5  & 14.0\% & $1400\times$ \\
    10 & 13.0\% & $1300\times$ \\
    15 &  8.7\% &  $870\times$ \\
    20 &  6.5\% &  $650\times$ \\
    \hline
  \end{tabular}
\end{table}

Content-Based KNN achieves \textbf{Precision@10 = 13\%} (1300$\times$ better than random),
produced without any user rating history, directly addressing the Cold Start problem.

\subsection{Evaluation --- TruncatedSVD Rating Prediction}

Figure~\ref{fig:svd_rmse} shows the TruncatedSVD evaluation on a random sample of 200 active users (80/20 split).

\begin{figure}[H]
  \centering
  \includegraphics[width=\linewidth]{3_svd_rmse_analysis.png}
  \caption{TruncatedSVD --- Actual vs.~Predicted ratings (left) and
           prediction error distribution (right).}
  \label{fig:svd_rmse}
\end{figure}

The RMSE of 2.5626 and MAE of 2.2801 are higher than the PyTorch SVD baseline.
This is because TruncatedSVD requires a dense matrix: unobserved entries are zero-filled
before decomposition, introducing a systematic downward bias in predicted scores.
This numerical bias does not invalidate the ranking: books predicted with higher
relative scores are still preferred by users, as confirmed qualitatively in
Figure~\ref{fig:sample_recs}.
For absolute rating prediction, the PyTorch SVD (Section~\ref{sec:svd_pytorch})
remains more accurate; TruncatedSVD is used here specifically for its ability to
efficiently process the full 53,424-user matrix for collaborative ranking signals.

\subsection{Qualitative Evaluation --- Sample Recommendations}

Figure~\ref{fig:sample_recs} compares recommendations by Content-Based KNN (Cold Start)
against the Hybrid system for an active user.

\begin{figure}[H]
  \centering
  \includegraphics[width=\linewidth]{4_sample_recommendations.png}
  \caption{Sample Recommendations: Content-Based KNN (left) vs.
           Hybrid KNN+SVD for active user~12874 (right).}
  \label{fig:sample_recs}
\end{figure}

For ``The Hunger Games'' (Cold Start), KNN correctly surfaces
\textit{Mockingjay} (\#1), \textit{Catching Fire} (\#2), and \textit{Harry Potter} (\#3)
using only content similarity.
For ``Harry Potter'' (Active User), the Hybrid system diversifies results
by blending collaborative signals, surfacing cross-genre titles that reflect the
specific user's taste profile rather than pure content similarity.

\subsection{Comprehensive Comparison: CF Baseline vs.~Hybrid Extension}

Table~\ref{tab:full_compare} provides a direct comparison of all models.

\begin{table}[H]
  \centering
  \caption{All Models Comparison --- CF Baseline vs. Hybrid Extension}
  \label{tab:full_compare}
  \begin{tabular}{|l|c|c|c|c|}
    \hline
    \textbf{Model} & \textbf{Cold Start?} & \textbf{Precision@10} & \textbf{RMSE} & \textbf{Discovery?} \\
    \hline
    User-Based KNN (CF)     & \textbf{Fails} & 1.3\% (CF-LOO)  & 1.0446 & No \\
    \hline
    Item-Based KNN (CF)     & \textbf{Fails} & 2.0\% (CF-LOO)  & 0.9811 & No \\
    \hline
    SVD-PyTorch (CF)        & \textbf{Fails} & 0.1\% (CF-LOO)  & 1.1537 & No \\
    \hline
    \textbf{CB-KNN (Hybrid)} & \textbf{Works} & \textbf{13\% (CB-LOO)} & N/A & \textbf{Yes} \\
    \hline
    \textbf{Hybrid KNN+SVD} & \textbf{Works} & \textbf{13\%+} (CB-LOO) & 2.5626* & \textbf{Yes} \\
    \hline
    \multicolumn{5}{l}{\small * TruncatedSVD RMSE affected by zero-fill bias; see Section~\ref{sec:hybrid} for explanation.} \\
    \multicolumn{5}{l}{\small CF-LOO and CB-LOO use different evaluation protocols; see Section~\ref{sec:hybrid} for details.} \\
  \end{tabular}
\end{table}

The key improvements of the Hybrid Extension are: (1) Cold Start is fully resolved ---
any user receives recommendations immediately without rating history;
(2) Discovery Mode supports free-text semantic queries absent in pure CF;
(3) Precision@10 improves substantially under the content-based evaluation protocol,
though this comparison should be read alongside the footnote regarding protocol differences.

% ─────────────────────────────────────────────────────────────────────────────
%  SECTION 6 – DISCUSSION
% ─────────────────────────────────────────────────────────────────────────────
\newpage
\section{Discussion}
\label{sec:discussion}

\subsection{Strengths and Weaknesses}

Table~\ref{tab:sw} summarises the key trade-offs of each method.

\begin{table}[H]
  \centering
  \caption{Strengths and weaknesses of each method}
  \label{tab:sw}
  \begin{tabular}{|p{2.5cm}|p{5.2cm}|p{5.2cm}|}
    \hline
    \textbf{Method} & \textbf{Strengths} & \textbf{Weaknesses} \\
    \hline
    User-Based KNN
      & No training; easy to update; high recommendation diversity
      & $O(n^2)$ with users; sensitive to sparse overlap; Cold Start \\
    \hline
    Item-Based KNN
      & Best on all 5 CF metrics; stable item profiles; explainable
      & Lower diversity; large item-item matrix; Cold Start for new books \\
    \hline
    SVD (PyTorch)
      & Handles sparsity via latent factors; explicit bias modelling; GPU-scalable
      & Low HitRate@10 (optimised for RMSE); requires retraining; less interpretable \\
    \hline
    Hybrid (CB-KNN + TruncatedSVD)
      & Solves Cold Start; supports free-text discovery; personalised for active users
      & TruncatedSVD RMSE bias from zero-fill; KNN brute-force $O(n \times d)$ per query \\
    \hline
  \end{tabular}
\end{table}

\subsection{Limitations}

\begin{table}[H]
  \centering
  \caption{Remaining limitations}
  \label{tab:limitations}
  \begin{tabular}{|p{2.5cm}|p{4.8cm}|p{4.5cm}|}
    \hline
    \textbf{Issue} & \textbf{CF Baseline} & \textbf{Hybrid Extension} \\
    \hline
    Cold Start & Completely fails & \textbf{Solved} by Content-Based KNN \\
    \hline
    Sparsity & Many zero-similarity scores & Partially mitigated: TF-IDF features are dense \\
    \hline
    SVD RMSE bias & N/A & Zero-fill bias in TruncatedSVD; rankings still accurate \\
    \hline
    Scalability & $O(n^2)$ similarity matrix & KNN brute-force: $O(n \times d)$ per query \\
    \hline
    Implicit feedback & Not supported & Not supported; ratings only \\
    \hline
  \end{tabular}
\end{table}

% ─────────────────────────────────────────────────────────────────────────────
%  SECTION 7 – CONCLUSION AND FUTURE WORK
% ─────────────────────────────────────────────────────────────────────────────
\newpage
\section{Conclusion and Future Work}
\label{sec:conclusion}

\subsection{Conclusion}

This project was carried out in two parts using the GoodBooks-10K dataset.

\textbf{Part~1 (CF Baseline)} built and evaluated three pure Collaborative Filtering models
under a Leave-One-Out protocol (8,000 users, 9,991 books):
Item-Based KNN achieved the best results across all five metrics (HitRate@10 = 20\%, RMSE = 0.9811);
User-Based KNN delivered solid diversity-oriented results (HitRate@10 = 12.8\%);
and the PyTorch SVD showed low HitRate@10 (1.4\%) due to its rating-prediction objective,
but has strong potential with BPR loss and extended training.
All three CF models fail on Cold Start.

\textbf{Part~2 (Hybrid Extension)} addressed the Cold Start problem:
Content-Based KNN (TF-IDF, 5,003 dimensions) achieves Precision@10 = 13\%
under the content-based LOO protocol, without requiring any user history.
Hybrid KNN+TruncatedSVD blends content and collaborative signals for active users,
producing diverse and personalised recommendations.
Smart Intent Detection enables a seamless single-input interface for
both exact title lookup and free-text semantic discovery.
The full system is deployed as a Flask web application.

\subsection{Future Work}

\begin{table}[H]
  \centering
  \caption{Proposed future directions}
  \label{tab:future}
  \begin{tabular}{|p{4cm}|p{9cm}|}
    \hline
    \textbf{Direction}       & \textbf{Technique} \\
    \hline
    Improve SVD ranking      & BPR Loss instead of MSE for ranking-aware training \\
    \hline
    Deeper architecture      & Neural CF, AutoRec, BERT4Rec for sequence-aware recommendations \\
    \hline
    Improve scalability      & Approximate Nearest Neighbours (FAISS, Annoy, HNSW) \\
    \hline
    Implicit feedback        & Matrix Factorization for implicit signals (ALS, BPR) \\
    \hline
    Fix SVD zero-fill bias   & Use mean-fill or matrix completion; train surprise SVD \\
    \hline
    Richer content features  & Book cover image embeddings (CLIP) for visual CB filtering \\
    \hline
  \end{tabular}
\end{table}

% ─────────────────────────────────────────────────────────────────────────────
%  REFERENCES
% ─────────────────────────────────────────────────────────────────────────────
\newpage
\begin{thebibliography}{99}

\bibitem{resnick1994}
P. Resnick, N. Iacovou, M. Suchak, P. Bergstrom, and J. Riedl,
``GroupLens: An open architecture for collaborative filtering of netnews,''
in \textit{Proc. ACM Conf. Computer Supported Cooperative Work (CSCW)},
pp. 175--186, 1994.

\bibitem{ricci2011}
F. Ricci, L. Rokach, and B. Shapira,
``Introduction to recommender systems handbook,''
in \textit{Recommender Systems Handbook}, Springer, pp. 1--35, 2011.

\bibitem{koren2009}
Y. Koren, R. Bell, and C. Volinsky,
``Matrix factorization techniques for recommender systems,''
\textit{Computer}, vol. 42, no. 8, pp. 30--37, Aug. 2009.

\bibitem{funk2006}
S. Funk,
``Netflix update: Try this at home,''
\textit{[Online]}. Available: \url{https://sifter.org/simon/journal/20061211.html}, Dec. 2006.

\bibitem{cover1967}
T. M. Cover and P. E. Hart,
``Nearest neighbor pattern classification,''
\textit{IEEE Trans. Inf. Theory}, vol. 13, no. 1, pp. 21--27, Jan. 1967.

\bibitem{su2009}
X. Su and T. M. Khoshgoftaar,
``A survey of collaborative filtering techniques,''
\textit{Advances in Artificial Intelligence}, vol. 2009, Article ID 421425, 2009.

\bibitem{schein2002}
A. I. Schein, A. Popescul, L. H. Ungar, and D. M. Pennock,
``Methods and metrics for cold-start recommendations,''
in \textit{Proc. 25th Annual Int. ACM SIGIR Conf. Research and Development in Information Retrieval},
pp. 253--260, 2002.

\bibitem{ziegler2005}
C.-N. Ziegler, S. M. McNee, J. A. Konstan, and G. Lausen,
``Improving recommendation lists through topic diversification,''
in \textit{Proc. 14th Int. Conf. World Wide Web (WWW)}, pp. 22--32, 2005.

\bibitem{bm25tfidf}
K. S. Jones,
``A statistical interpretation of term specificity and its application in retrieval,''
\textit{Journal of Documentation}, vol. 28, no. 1, pp. 11--21, 1972.

\bibitem{sklearn}
F. Pedregosa et al.,
``Scikit-learn: Machine learning in Python,''
\textit{Journal of Machine Learning Research}, vol. 12, pp. 2825--2830, 2011.

\end{thebibliography}

\end{document}